\par\noindent\textbf{(1)} 小位移w满足 μ\_s w\_tt=T\_s(w\_xx+w\_yy)。固定端本征模 φ\_mn=sin(mπx/a)sin(nπy/b)，m,n≥1；代入得 ω\_mn²=(T\_s/μ\_s)π²(m²/a²+n²/b²)。模态正交内积为∫φ\_mnφ\_pq dA∝δ\_mpδ\_nq。
\par\noindent\textbf{(2)} 若a=b，频率只依赖m²+n²，因此m≠n时φ\_mn和φ\_nm通常简并。该简并源于正方形对称群；即使两个整数对偶然有同一平方和，也应按整个简并子空间处理，不可以普通单模微扰直接分配频移。
\par\noindent\textbf{(3)} 由广义本征方程 −T\_s∇²φ=ω² μ\_sφ，对面密度扰动求一阶Rayleigh商，得 δω/ω=−(1/2)[∫δμ|φ|²dA]/[μ\_s∫|φ|²dA]。局部加质量对正质量扰动必降低频率，最大下降发生在未扰动模态的反节点位置。
\par\noindent\textbf{(4)} 简并模φ\_1、φ\_2需对矩阵 K\_ij=−(ω₀/2μ\_s)∫δμ φ\_iφ\_j dA/∫φ\_i²dA 作Hermitian对角化（基底先归一化）。若扰动保留交换 x↔y 对称性，两种对称/反对称组合可成为新模；若矩阵正比单位阵则一阶不裂分。
